% Review-only v10 cumulative TODO overlay, 9 October 2026.
% To be consolidated with the previous Lemma 6.49 review note, NOT
% simultaneously inserted as a second duplicate note.
% The author's four separate reconstructed mathematical repairs are untouched.
%
% Suggested precise replacement at Lemma 6.49:
\todo[inline]{\textbf{v10 validation (actual first full type difference).}
The common-ambient argument now includes the complete finite
$L^+$-prefix records, not just the directed binary trace.
For two vertices of $C^+$ having the same root type, the
unique largest common induced partial-type prefix below any
finite bound can be selected canonically. If this bound is
below both $E$-free levels and the common prefix ends before
the bound, the next coordinate must exhibit a genuine
difference in one orientation of a directed binary $L$-relation.
All unary, diagonal and auxiliary $E$-coordinates have
been accounted for. See \texttt{FirstFullPrefix.lean} and
\texttt{CommonSocleType.lean} (after checking the exact commit).
To identify the resulting coordinate with the tree meet,
one must still identify induced restrictions from a common
ambient $C^+$ with the predecessors in the *actual*
$K^\mathcal F_{pt}$ tree and verify the proposed charged
coordinate lies below both free levels. In the positive
non-top mixed-generation case, the required guard is the
strict $p+1<j$, not merely $p<j$.}
%
% Proof endpoints, conditional on green exact-head CI:
%   fullPartialType_eq_at_smaller
%   exists_maximal_common_fullPrefix
%   maximal_common_fullPrefix_unique
%   binary_witness_of_first_fullPrefix_difference
%   exists_maximal_common_fullPrefix_with_binary_witness
